% exercises_preamble.tex
%
% Präambel für die Übungsblätter zum Selbststudium (Kapitel N.3 bzw. N.5) als
% PDF, übernommen aus der Präambel der CA-2-Unterlagen
% (internal/exams/exam_2026ws/ca02/preamble.tex). Die Studierenden üben damit
% das Format von CA 3: Aufgaben auf Papier, gelöst lokal am PC.
% Die .tex-Dateien erzeugt convert_myst_to_latex.py aus dem Master.

% ── Seitenlayout ─────────────────────────────────────────────────────────
\usepackage[a4paper, width=150mm, height=240mm]{geometry}

\usepackage{amsmath, amssymb}

\usepackage[no-math]{fontspec}
\usepackage{unicode-math}
\usepackage{libertinus}

\usepackage{setspace}
\usepackage{microtype}
\usepackage{textcomp}
\usepackage[main=english]{babel}

% ✩ für den Schwierigkeitsgrad (Libertinus hat das Zeichen nicht)
\usepackage{pifont}
\usepackage{newunicodechar}
\newunicodechar{✩}{\ding{73}}

% ── Kopf-/Fußzeilen ──────────────────────────────────────────────────────
\usepackage[automark]{scrlayer-scrpage}
\clearpairofpagestyles

\ifoot{S. Gramsch}
\cfoot{TH Mannheim -- Applied Numerical Mathematics -- WS 2026/27}
\ofoot{\pagemark}

% Keine Titelei: Die erzeugte .tex-Datei beginnt mit der Überschrift
% "N.M Exercises" als \section*, darunter "How to Practice" und "Exercises"
% als \subsection*.

% ── Listen ───────────────────────────────────────────────────────────────
\usepackage{tabularx}
\usepackage{booktabs}
\usepackage{enumitem}

% ── Farben ───────────────────────────────────────────────────────────────
\usepackage{xcolor}
\definecolor{itwm_blue_04}{HTML}{005A94}
\definecolor{frankfurtBlueLight}{HTML}{AECAE2}

\definecolor{my_darkgray}{HTML}{484949}
\definecolor{my_lightgray}{HTML}{F3F4F4}
\definecolor{my_orange}{HTML}{E87846}

% Farben der Admonition "tip" im Buch (MyST book-theme, heller Modus):
% Rand und Symbol --myst-color-success, Kopfzeile --myst-color-success-bg
\definecolor{tip_green}{HTML}{16A34A}
\definecolor{tip_bg}{HTML}{F0FDF4}

% ── Python-Code ──────────────────────────────────────────────────────────
\usepackage{listings}

% Zeilenumbrüche im Code mit Pfeil markieren, damit beim Abtippen klar ist,
% dass es eine einzige Python-Zeile ist
\lstdefinestyle{python}{
  language=Python,
  basicstyle=\ttfamily\small\setstretch{1},
  keywordstyle=\color{itwm_blue_04}\bfseries,
  stringstyle=\color{my_orange!80!black},
  commentstyle=\color{gray},
  showstringspaces=false,
  columns=fullflexible,
  keepspaces=true,
  upquote=true,
  breaklines=true,
  postbreak=\mbox{\textcolor{gray}{$\hookrightarrow$}\space},
  extendedchars=true,
}

\lstdefinestyle{ausgabe}{
  basicstyle=\ttfamily\footnotesize\setstretch{1},
  columns=fullflexible,
  keepspaces=true,
  upquote=true,
}

\lstdefinestyle{inline}{
  basicstyle=\ttfamily,
  showstringspaces=false,
  columns=fullflexible,
  keepspaces=true,
  upquote=true,
}

% Code im Fließtext: \py|derivative(values, dt)|
\newcommand{\py}{\lstinline[style=inline]}

% ── Kästen ───────────────────────────────────────────────────────────────
\usepackage[most]{tcolorbox}
\tcbuselibrary{listings}

%% Aufgabe ---------------------------------------------------------------
% \begin{aufgabe}{Exercise 7.1 (✩)}{approx.\ 7 min}: Bezeichnung links,
% Richtzeit rechts (in CA 2 und CA 3 stehen dort die Punkte).
% Aussehen wie die Admonition "tip" im Buch: grüner Rand links, hellgrüne
% Kopfzeile mit grünem Stift, weißer Inhalt, leichter Schatten.
% Eine Aufgabe wird nicht über zwei Seiten geteilt: Passt sie nicht mehr,
% beginnt sie auf der nächsten Seite. Nur für Aufgaben, die länger als eine
% Seite sind, setzt der Konverter die teilbare Variante aufgabelang ein.
\tcbset{
  aufgabenstil/.style 2 args={
    enhanced,
    boxrule=0pt,
    frame hidden,
    borderline west={1.5pt}{0pt}{tip_green},
    title={\textcolor{tip_green}{\ding{46}}\enspace #1 \hfill #2},
    coltitle=black,
    fonttitle=\bfseries,
    colbacktitle=tip_bg,
    colback=white,
    drop fuzzy shadow=black!25,
    arc=1mm,
    left=4mm, right=4mm,
    before skip=4mm, after skip=4mm,
  },
}
\newtcolorbox{aufgabe}[2]{aufgabenstil={#1}{#2}}
\newtcolorbox{aufgabelang}[2]{aufgabenstil={#1}{#2}, breakable}

%% Code und Ausgabe ------------------------------------------------------
% breakable: längerer vorgegebener Code darf über eine Seite hinausgehen,
% sonst bleibt vor ihm eine halbe Seite leer
\newtcblisting{code}{
    listing only,
    enhanced,
    breakable,
    boxrule=0pt,
    frame hidden,
    sharp corners,
    colback=my_lightgray,
    left=4pt, right=4pt, top=2pt, bottom=2pt,
    listing options={style=python},
}

\newtcblisting{ausgabe}{
    listing only,
    enhanced,
    boxrule=0pt,
    frame hidden,
    sharp corners,
    colback=white,
    borderline west={2pt}{0pt}{my_lightgray!70!black},
    left=6pt, right=4pt, top=0pt, bottom=0pt,
    listing options={style=ausgabe},
}

% ── Verweise ─────────────────────────────────────────────────────────────
\usepackage[hidelinks]{hyperref}
